\section*{Capítulo \ref{ch:g11:subduction-zones} --- As zonas de subducção}

\begin{solution}{exo:g11:subduction-zones:1}
A \omterm{def:g11:subduction-zones:subduction}{subducção} é o afundamento de uma placa de \omterm{def:g11:ocean-ridges:oceanic-lithosphere}{litosfera oceânica} sob outra
placa, de volta para o manto; uma \omterm{def:g11:subduction-zones:subduction}{zona de subducção} é um limite de placas
convergente onde ela acontece, marcado por uma fossa. É a \omterm{def:g11:ocean-ridges:oceanic-lithosphere}{litosfera oceânica}
que afunda: quando é antiga, ela é fria, espessa e mais densa que a
\omterm{def:g11:lithosphere-asthenosphere:lithosphere}{astenosfera} abaixo dela. A crosta continental é leve demais para afundar.
\end{solution}

\begin{solution}{exo:g11:subduction-zones:2}
(1) Por exemplo a região de Tonga--Fiji, o Pacífico ocidental do Japão às
Marianas, a Indonésia e a América do Sul. (2) Sempre do lado da placa que se
sobrepõe, nunca do lado do oceano. (3) Rasos: todos os terremotos da Dorsal
Mesoatlântica no mapa têm menos de \qty{70}{km} de profundidade.
\end{solution}

\begin{solution}{exo:g11:subduction-zones:3}
$7031 \times 0.21 \approx 1480$ terremotos com mais de \qty{70}{km} de
profundidade; $7031 \times 0.069 \approx 485$ com mais de \qty{300}{km}.
\end{solution}

\begin{solution}{exo:g11:subduction-zones:4}
O \omterm{def:g11:subduction-zones:eclogite}{xisto azul} é uma rocha metamórfica formada a partir do basalto da \omterm{def:g11:ocean-ridges:oceanic-lithosphere}{crosta
oceânica} enterrado sob alta pressão enquanto continua relativamente frio,
colorido de azul pela glaucofana; o \omterm{def:g11:subduction-zones:eclogite}{eclogito} é a rocha metamórfica densa em
que o basalto e o \omterm{def:g11:ocean-ridges:gabbro}{gabro} se transformam em profundidades do manto, feita de
granada vermelha e de piroxênio verde. Os dois vêm da \omterm{def:g11:ocean-ridges:oceanic-lithosphere}{crosta oceânica} e
conservam a sua composição química; o que difere são os minerais, formados
no estado sólido sob a nova pressão e a nova temperatura, e a densidade.
\end{solution}

\begin{solution}{exo:g11:subduction-zones:5}
Fundo oceânico: \omterm{def:g11:lithosphere-asthenosphere:lithosphere}{litosfera} antiga que se move em direção à fossa. Fossa: a
placa se dobra para baixo e começa a afundar. Antearco: os sedimentos
raspados da placa que afunda se empilham; abaixo, as duas placas deslizam
uma contra a outra no limite de placas, onde acontecem grandes terremotos.
Arco vulcânico: uma linha de vulcões acima da \omterm{def:g11:subduction-zones:slab}{laje}. Retroarco: atrás de
alguns arcos, a placa que se sobrepõe é estirada.
\end{solution}

\begin{solution}{exo:g11:subduction-zones:6}
$\tan\delta = (508 - 135)/(800 - 300) = 373/500 \approx 0.746$, logo
$\delta \approx \ang{37}$.
\end{solution}

\begin{solution}{exo:g11:subduction-zones:7}
(1) $\tan\ang{45} = 1$: \qty{300}{km} de profundidade. (2) A \qty{670}{km}
da fossa. (3) $670\sqrt{2} \approx \qty{948}{km}$ de \omterm{def:g11:subduction-zones:slab}{laje}.
\end{solution}

\begin{solution}{exo:g11:subduction-zones:8}
(1) $(74.6 + 70.0 + 65.5 + 63.1 + 52.6)/5 \approx \qty{65.2}{mm}$ por ano;
$(47.6 + 31.0 + 29.2 + 22.8 + 19.8 + 15.2 + 14.9 + 10.6)/8 \approx
\qty{23.9}{mm}$ por ano; razão $65.2/23.9 \approx 2.7$. (2) As placas se
movem mais depressa quando uma \omterm{def:g11:subduction-zones:slab}{laje} pende delas: a \omterm{def:g11:subduction-zones:slab-pull}{tração da laje} que afunda
é a força mais forte que move as placas, como concluíram Forsyth e Uyeda.
\end{solution}

\begin{solution}{exo:g11:subduction-zones:9}
(1) \qty{83}{mm} equivalem a \qty{0.083}{m}: $50/0.083 \approx 600$ anos e
$60/0.083 \approx 720$ anos de movimento da placa. (2) De 869 a 2011 são
1142 anos: da mesma ordem. Se o limite de placas fica travado entre os
grandes terremotos, o deslizamento de 2011 liberou vários séculos de
movimento acumulado.
\end{solution}

\begin{solution}{exo:g11:subduction-zones:10}
(1) Japão: $6400 - 3200\,e^{-133/62.8} \approx \qty{6015}{m}$; Tonga:
$6400 - 3200\,e^{-99/62.8} \approx \qty{5740}{m}$; Andes:
$2500 + 350\sqrt{49} = \qty{4950}{m}$. (2) O fundo que chega ao Japão: o
mais antigo, portanto a \omterm{def:g11:lithosphere-asthenosphere:lithosphere}{litosfera} mais fria e mais densa, e o fundo mais
profundo.
\end{solution}

\begin{solution}{exo:g11:subduction-zones:11}
(1) $3380 \times 1.06 \approx \qty{3583}{kg/m^3}$. (2)
$(3583 - 2900)/2900 \approx 0.24$: um aumento de cerca de \qty{24}{\percent}.
(3) Antes, a crosta é mais leve que o manto em
$3380 - 2900 = \qty{480}{kg/m^3}$; depois, ela é mais densa em cerca de
\qty{203}{kg/m^3}. Ela já não sustenta a \omterm{def:g11:subduction-zones:slab}{laje}: ela a puxa para baixo.
\end{solution}

\begin{solution}{exo:g11:subduction-zones:12}
(1) $t = 6 \times 10^6 \times 3.16 \times 10^7 \approx 1.9 \times 10^{14}$~s;
$\sqrt{10^{-6} \times 1.9 \times 10^{14}} \approx \qty{1.4e4}{m}$, cerca de
\qty{14}{km}. (2) Muito menos que \qty{100}{km}: em 6 milhões de anos, só a
pele externa da \omterm{def:g11:subduction-zones:slab}{laje} se aquece; o seu interior continua frio. (3) A rocha
fria é rígida e pode se romper, enquanto o manto quente em volta escoa; os
terremotos profundos acontecem dentro da \omterm{def:g11:subduction-zones:slab}{laje} fria.
\end{solution}

\begin{solution}{exo:g11:subduction-zones:13}
(1) $\sqrt{1017^2 + 670^2} \approx \qty{1218}{km}$. (2) \qty{64}{mm} por ano
equivalem a \qty{64}{km} por milhão de anos: $1218/64 \approx 19$ milhões de
anos. (3) Cerca de três vezes mais tempo que sob Tonga: a \omterm{def:g11:subduction-zones:slab}{laje} sob os Andes
afunda mais devagar e com um mergulho mais suave.
\end{solution}

\begin{solution}{exo:g11:subduction-zones:14}
(1) Uma fatia de \qty{1}{m} de largura ao longo da fossa tem um volume de
$10^5 \times 10^6 \times 1 = \qty{1e11}{m^3}$; o seu excesso de massa é
$36 \times 10^{11} = \qty{3.6e12}{kg}$. (2) O seu excesso de peso é
$3.6 \times 10^{12} \times 9.8 \approx \qty{3.5e13}{N}$ para cada metro de
fossa. (3) A \omterm{def:g11:subduction-zones:slab}{laje} está presa ao resto da placa, que é rígida: a força é
transmitida à placa inteira, que é arrastada em direção à fossa, ainda mais
porque as fossas se estendem por milhares de quilômetros.
\end{solution}

\begin{solution}{exo:g11:subduction-zones:15}
(1) $(23.6 - 16.5) \times 111 \approx \qty{790}{km}$. (2) A taxa cresce
$237 - 104 = \qty{133}{mm}$ por ano em \qty{790}{km}, cerca de \qty{17}{mm}
por ano a cada cem quilômetros. (3) A convergência é a velocidade da placa
do Pacífico menos a do arco. Como a velocidade da placa do Pacífico é mais
ou menos a mesma nos dois pontos, o próprio arco deve estar se movendo em
direção à placa do Pacífico, cada vez mais depressa para o norte: o
retroarco atrás dele está sendo estirado, mais no norte do que no sul.
\end{solution}

\begin{solution}{pb:g11:subduction-zones:1}
\textbf{1.} Cerca de \qty{671}{km}.

\textbf{2.} Do lado do oceano, a placa ainda não afundou: a \omterm{def:g11:lithosphere-asthenosphere:lithosphere}{litosfera} ali
fica perto da superfície. Os terremotos profundos acontecem dentro da \omterm{def:g11:subduction-zones:slab}{laje},
que fica apenas sob a placa que se sobrepõe.

\textbf{3.} $z(300) = -24 + 273 = \qty{249}{km}$ e
$z(700) = -24 + 637 = \qty{613}{km}$; $\tan\delta = 364/400 = 0.91$, logo
$\delta \approx \ang{42}$.

\textbf{4.} $x = (100 + 24)/0.91 \approx \qty{136}{km}$ e
$x = (400 + 24)/0.91 \approx \qty{466}{km}$.

\textbf{5.} Uma reta com esse mergulho que começasse na fossa passaria por
$z = 0$ em $x = 0$; a reta fica acima. A \omterm{def:g11:subduction-zones:slab}{laje} não é reta: perto da fossa,
ela mergulha mais suavemente, e fica mais íngreme mais abaixo.

\textbf{6.} \qty{16.8}{cm} por ano, \qty{168}{km} por milhão de anos.

\textbf{7.} $168/69 \approx 2.4$ vezes mais rápida.

\textbf{8.} A convergência através da fossa é a velocidade da placa do
Pacífico menos a velocidade do arco do outro lado. Se o arco se move em
direção à placa do Pacífico, como acontece enquanto o retroarco atrás dele é
estirado, as duas velocidades se somam, e a convergência é maior que a
velocidade da placa do Pacífico sozinha.

\textbf{9.} $6400 - 3200\,e^{-99/62.8} \approx \qty{5740}{m}$.

\textbf{10.} Depois de cerca de 60 milhões de anos, a \omterm{def:g11:ocean-ridges:oceanic-lithosphere}{litosfera oceânica} é,
como um todo, mais densa que a \omterm{def:g11:lithosphere-asthenosphere:lithosphere}{astenosfera} abaixo dela; aos 99 milhões de
anos, ela é ainda mais densa. Assim que a sua borda desce, o seu peso puxa o
resto para dentro.

\textbf{11.} $(670 + 24)/0.91 \approx \qty{763}{km}$.

\textbf{12.} $\sqrt{763^2 + 670^2} \approx \qty{1015}{km}$.

\textbf{13.} $1015/168 \approx 6.0$ milhões de anos.

\textbf{14.} $1015/69 \approx 14.7$ milhões de anos.

\textbf{15.} $3380 \times 1.06 \approx \qty{3583}{kg/m^3}$: a crosta, antes
mais leve que o manto, fica mais densa que ele, e a \omterm{def:g11:subduction-zones:slab}{laje} inteira puxa com
mais força.

\textbf{16.} $t \approx 6.0 \times 10^6 \times 3.16 \times 10^7 \approx
1.9 \times 10^{14}$~s, e $\sqrt{\kappa t} \approx \qty{14}{km}$.

\textbf{17.} \qty{14}{km} é pouco ao lado de \qty{100}{km}: só a parte
externa da \omterm{def:g11:subduction-zones:slab}{laje} se aqueceu; o seu interior ainda está frio quando ela chega
a \qty{670}{km}.

\textbf{18.} A \omterm{def:g11:subduction-zones:slab}{laje} fria continua rígida o bastante para se romper, até a
sua ponta, enquanto o manto quente em volta escoa: os terremotos sob Tonga
acontecem dentro da \omterm{def:g11:subduction-zones:slab}{laje}, até cerca de \qty{671}{km}.

\textbf{19.} As ondas sísmicas mostram uma mudança brusca no manto a cerca
de \qty{670}{km}: as suas velocidades dão um salto ali.

\textbf{20.} Cerca de \textbf{6 milhões de anos} da fossa até
\qty{670}{km}.
\end{solution}
